\slajdid{09_1-s035}
\begin{frame}[label=09_1-s035]
\gusttekst\alert{Invertor u nMOS tehnologiji sa aktivnim opterećenjem na izlazu\par u vidu nMOS tranzistora sa indukovanim kanalom}\par\medskip
Kod invertora u nMOS tehnologiji sa pasivnim opterećenjem na izlazu nismo prodiskutovali izbor pasivnog opterećenja, odnosno otpornika $R_L$.
\begin{enumerate}
\item Otpornik $R_L$ treba da je što veći kako bi napon logičke nule bio što manji. Otpornik $R_L$ treba da bude što je veći, da bi što manje uticao na strujni kapacitet logičke nule.
\item Otpornik $R_L$ treba da je što manji, da bi što brže punio kapacitivnosti na izlazu.
\item Otpornik $R_L$ direktno utiče na $V_{IL}$ i $V_{IH}$. Što je manji $V_{IL}$ je veće, ali je $V_{IH}$ veće i obrnuto.
\end{enumerate}
Znači potreban je kompromis za izbor vrednosti pasivnog opterećenja. Uočite da se praktično u svim izrazima pojavljuje faktor $k_nR_L$, odnosno postavlja se pitanje izbora tog faktora. A na njega možemo kod MOS tranzistora da utičemo promenom faktora $k_n$. Da se podsetimo
\[k_n=\mu_nC_{oxn}\frac{\alert{W_n}}{L_n}\]
i ako su parametri $\mu_n$, $C_{oxn}$ i $L_n$ određeni tehnološkim mogućnostima i fiksni (želimo na primer tranzistore minimalnih dimenzija pa je $L_n=L_{nmin}$ koje tehnologija dopušta) onda promenom parametra $W_n$ na lak način možemo da utičemo na $k_n$. \alert{$W_n$ je širina kanala i nju u dizajnu možemo lako da kontrolišemo.}
\end{frame}
